\documentclass[11pt,a4paper]{article}
\usepackage[utf8]{inputenc}
\usepackage[spanish,es-tabla]{babel}
\usepackage{amsmath,amsfonts,amssymb}
\usepackage{geometry}
\usepackage{graphicx}
\usepackage{xcolor}

\usepackage{enumitem}

\geometry{margin=2.5cm}

\title{\textbf{Soluci\'on PC Not\'ee 2021-2022\\TAF ISC - UEC 2\\28/10/2021}}
\author{}
\date{}

\begin{document}

\maketitle

\section*{Sistema de modulaci\'on 8-PSK}

\textbf{Datos del problema:}
\begin{itemize}
    \item Modulaci\'on: 8-PSK
    \item D\'ebito binario: $D_b = 48$ Mbps
    \item Frecuencia portadora: $\nu_0 = 2.4$ GHz
    \item Elementos binarios independientes e id\'enticamente distribuidos con probabilidad igual en $\{0,1\}$
\end{itemize}

\textbf{Se\~nal transmitida:}
\[
s(t) = A\sum_k \left[a_k h(t-kT)\cos(2\pi\nu_0 t) - b_k h(t-kT)\sin(2\pi\nu_0 t)\right]
\]

\textbf{Funci\'on de transferencia del filtro:}
\[
H(\nu) = \begin{cases}
\sqrt{T} & |\nu| \leq \frac{1}{3T} \\
\sqrt{2T\left(1 - \frac{3T}{2}|\nu|\right)} & \frac{1}{3T} \leq |\nu| \leq \frac{2}{3T} \\
0 & \text{en otro caso}
\end{cases}
\]

\textbf{Canal:} AWGN con ruido $n(t)$ de media cero y densidad espectral de potencia $\frac{N_0}{2}$

\textbf{Se\~nal recibida:} $r(t) = s(t) + n(t)$

\subsection*{Pregunta 1: Alfabeto 8-PSK}

\textbf{Enunciado:} Definimos $d_k = a_k + ib_k$ y $\mathcal{D}$ el conjunto de todos los valores posibles de $d_k$. Dar el alfabeto 8-PSK $\mathcal{D}$.

\vspace{0.3cm}
\textbf{Soluci\'on:}

El alfabeto 8-PSK consiste en 8 s\'imbolos uniformemente distribuidos en un c\'irculo:

\[
\boxed{\mathcal{D} = \left\{e^{i\frac{(2k+1)\pi}{8}} \mid k \in \{0, 1, 2, 3, 4, 5, 6, 7\}\right\}}
\]

Expresado en forma expl\'icita:
\[
\mathcal{D} = \left\{e^{i\frac{\pi}{8}}, e^{i\frac{3\pi}{8}}, e^{i\frac{5\pi}{8}}, e^{i\frac{7\pi}{8}}, e^{i\frac{9\pi}{8}}, e^{i\frac{11\pi}{8}}, e^{i\frac{13\pi}{8}}, e^{i\frac{15\pi}{8}}\right\}
\]

\textbf{Representaci\'on gr\'afica:} Los 8 puntos est\'an en los \'angulos $\frac{\pi}{8}, \frac{3\pi}{8}, \frac{5\pi}{8}, \ldots, \frac{15\pi}{8}$, todos con m\'odulo unitario (si normalizamos a amplitud 1).

\textbf{Interpretaci\'on:}
\begin{itemize}
    \item Cada s\'imbolo codifica $\log_2(8) = 3$ bits
    \item Los s\'imbolos est\'an equiespaciados angularmente (cada $45°$)
    \item Todos tienen la misma energ\'ia (modulaci\'on de fase pura)
\end{itemize}



\subsection*{Pregunta 2: Mapping de Gray}

\textbf{Enunciado:} Proponer una conversi\'on binario a s\'imbolo que satisfaga la regla de Gray mapping.

\vspace{0.3cm}
\textbf{Soluci\'on:}

El Gray mapping asegura que s\'imbolos adyacentes difieran en solo un bit. Una propuesta es:

\begin{center}
\begin{tabular}{|c|c|c|}
\hline
\textbf{\'Angulo} & \textbf{S\'imbolo} & \textbf{C\'odigo Gray} \\
\hline
$\frac{\pi}{8}$ & $e^{i\pi/8}$ & 000 \\
$\frac{3\pi}{8}$ & $e^{i3\pi/8}$ & 001 \\
$\frac{5\pi}{8}$ & $e^{i5\pi/8}$ & 011 \\
$\frac{7\pi}{8}$ & $e^{i7\pi/8}$ & 010 \\
$\frac{9\pi}{8}$ & $e^{i9\pi/8}$ & 110 \\
$\frac{11\pi}{8}$ & $e^{i11\pi/8}$ & 111 \\
$\frac{13\pi}{8}$ & $e^{i13\pi/8}$ & 101 \\
$\frac{15\pi}{8}$ & $e^{i15\pi/8}$ & 100 \\
\hline
\end{tabular}
\end{center}

\textbf{Verificaci\'on del Gray mapping:}
\begin{itemize}
    \item $000 \rightarrow 001$: difieren en 1 bit ✓
    \item $001 \rightarrow 011$: difieren en 1 bit ✓
    \item $011 \rightarrow 010$: difieren en 1 bit ✓
    \item Cada par de s\'imbolos adyacentes difiere en exactamente 1 bit
\end{itemize}

\textbf{Ventaja:} Minimiza la tasa de error binaria cuando ocurren errores en s\'imbolos adyacentes.



\subsection*{Pregunta 3: Rapidez de s\'imbolo y valor de T}

\textbf{Enunciado:} Calcular la rapidez de s\'imbolo y deducir el valor de $T$.

\vspace{0.3cm}
\textbf{Soluci\'on:}

La relaci\'on entre d\'ebito binario y rapidez de s\'imbolo es:
\[
D_s = \frac{D_b}{\log_2(M)}
\]

Para 8-PSK: $M = 8$, por lo tanto:
\[
D_s = \frac{48 \times 10^6}{3} = \boxed{16 \times 10^6 \text{ s\'imbolos/s} = 16 \text{ Mbauds}}
\]

El periodo s\'imbolo es:
\[
T = \frac{1}{D_s} = \frac{1}{16 \times 10^6} = \boxed{6.25 \times 10^{-8} \text{ s} = 62.5 \text{ ns}}
\]

\textbf{Interpretaci\'on:}
\begin{itemize}
    \item Cada s\'imbolo transporta 3 bits
    \item Se transmiten 16 millones de s\'imbolos por segundo
    \item Cada s\'imbolo dura 62.5 nanosegundos
\end{itemize}



\subsection*{Pregunta 4: Banda ocupada y condici\'on narrowband}

\textbf{Enunciado:} Calcular la banda ocupada. ¿Se satisface la condici\'on de banda estrecha?

\vspace{0.3cm}
\textbf{Soluci\'on:}

\textbf{C\'alculo de la banda ocupada:}

Del soporte de $H(\nu)$, vemos que el filtro pasa-bajo tiene banda $\left[-\frac{2}{3T}, \frac{2}{3T}\right]$.

Para la se\~nal modulada, la banda ocupada es:
\[
B_{occ} = 2 \times \frac{2}{3T} = \frac{4}{3T}
\]

Sustituyendo $T = 6.25 \times 10^{-8}$:
\[
B_{occ} = \frac{4}{3 \times 6.25 \times 10^{-8}} = \frac{4}{1.875 \times 10^{-7}} = \boxed{21.33 \times 10^6 \text{ Hz} = 21.33 \text{ MHz}}
\]

\textbf{Verificaci\'on de la condici\'on narrowband:}

La condici\'on de banda estrecha requiere:
\[
B_{occ} \ll \nu_0
\]

Comparamos:
\begin{align*}
B_{occ} &= 21.33 \times 10^6 \text{ Hz} \\
\nu_0 &= 2.4 \times 10^9 \text{ Hz}
\end{align*}

Ratio:
\[
\frac{B_{occ}}{\nu_0} = \frac{21.33 \times 10^6}{2.4 \times 10^9} = 0.0089 \approx 0.9\%
\]

\textbf{Conclusi\'on:} $\boxed{\text{SÍ, la condici\'on narrowband se satisface}}$ ya que $B_{occ}$ es menor del 1\% de $\nu_0$.

\textbf{Consecuencia:} Podemos usar aproximaciones de banda estrecha en el an\'alisis del sistema.



\subsection*{Pregunta 5: Estructura del receptor ML \'optimo}

\textbf{Enunciado:} Dar la estructura del receptor \'optimo de m\'axima verosimilitud (ML). (No se demanda demostraci\'on).

\vspace{0.3cm}
\textbf{Soluci\'on:}

\textbf{Estructura del receptor \'optimo ML:}

\begin{center}
\begin{tabular}{c}
\texttt{r(t)} \\
$\downarrow$ \\
\hline
\textbf{Rama I (In-phase):} \\
\fbox{$\times \cos(2\pi\nu_0 t)$} $\rightarrow$ \fbox{Filtro PB $g(t) = h(t_0-t)$} $\rightarrow$ \fbox{Muestreo $t_0+kT$} $\rightarrow a_k'$ \\
\hline
\textbf{Rama Q (Quadrature):} \\
\fbox{$\times (-\sin(2\pi\nu_0 t))$} $\rightarrow$ \fbox{Filtro PB $g(t) = h(t_0-t)$} $\rightarrow$ \fbox{Muestreo $t_0+kT$} $\rightarrow b_k'$ \\
\hline
$\downarrow$ \\
\fbox{Formar $z_k = a_k' + ib_k'$} \\
$\downarrow$ \\
\fbox{Decisor ML: $\hat{d}_k = \arg\min_{d \in \mathcal{D}} |z_k - d|$} \\
$\downarrow$ \\
\texttt{Bits estimados}
\end{tabular}
\end{center}

\textbf{Descripci\'on de los bloques:}

\begin{enumerate}
    \item \textbf{Demoduladores I/Q:} Separan las componentes en fase y cuadratura multiplicando por $\cos(2\pi\nu_0 t)$ y $-\sin(2\pi\nu_0 t)$
    
    \item \textbf{Filtros adaptados:} $g(t) = h(t_0 - t)$ maximizan el SNR y eliminan componentes de alta frecuencia
    
    \item \textbf{Muestreadores:} Toman muestras en los instantes \'optimos $t_0 + kT$
    
    \item \textbf{Formador de s\'imbolo complejo:} Combina las salidas I/Q: $z_k = a_k' + ib_k'$
    
    \item \textbf{Decisor ML:} Elige el s\'imbolo m\'as cercano de $\mathcal{D}$:
    \[
    \hat{d}_k = \arg\min_{d \in \mathcal{D}} |z_k - d|^2
    \]
    
    \item \textbf{Demapping:} Convierte el s\'imbolo decidido a bits usando la tabla de Gray mapping inversa
\end{enumerate}

\textbf{Propiedades:}
\begin{itemize}
    \item Minimiza la probabilidad de error de s\'imbolo
    \item Maximiza el SNR
    \item Permite detecci\'on s\'imbolo por s\'imbolo si se cumple TSWOC
\end{itemize}



\subsection*{Pregunta 6: Detecci\'on s\'imbolo por s\'imbolo y TSWOC}

\textbf{Enunciado:} Mostrar que la detecci\'on \'optima puede hacerse s\'imbolo por s\'imbolo. Deducir que $\int_{\mathbb{R}} |h(t)|^2\, dt = 1$.

\vspace{0.3cm}
\textbf{Soluci\'on:}

\textbf{Demostraci\'on del TSWOC:}

El criterio TSWOC se satisface si:
\[
\sum_p |H(\nu - p/T)|^2 = T, \quad \forall \nu \in \mathbb{R}
\]

Analizamos la suma en diferentes intervalos:

\textbf{Intervalo 1:} $\left[\frac{p}{T} - \frac{1}{3T}, \frac{p}{T} + \frac{1}{3T}\right]$

Solo $|H(\nu - p/T)|^2$ es no nulo:
\[
\sum_k |H(\nu - k/T)|^2 = |H(\nu - p/T)|^2 = T \quad ✓
\]

\textbf{Intervalo 2:} $\left[\frac{p}{T} + \frac{1}{3T}, \frac{p}{T} + \frac{2}{3T}\right]$

Son no nulos $|H(\nu - p/T)|^2$ y $|H(\nu - (p+1)/T)|^2$:

\begin{align*}
\sum_k |H(\nu - k/T)|^2 &= |H(\nu - p/T)|^2 + |H(\nu - (p+1)/T)|^2 \\
&= 2T\left(1 - \frac{3T}{2}(\nu - p/T)\right) + 2T\left(1 - \frac{3T}{2}(-(nu - (p+1)/T))\right) \\
&= 2T\left(1 - \frac{3T}{2}(\nu - p/T)\right) + 2T\left(1 - \frac{3T}{2}(-\nu + p/T + 1/T)\right) \\
&= 2T + 2T\left(1 - \frac{3}{2}\right) \\
&= T \quad ✓
\end{align*}

\textbf{Conclusi\'on:} El TSWOC se satisface: $\sum_p |H(\nu - p/T)|^2 = T$ para todo $\nu$.

\textbf{Consecuencia - Detecci\'on s\'imbolo por s\'imbolo:}

Como TSWOC se satisface:
\[
\int_{\mathbb{R}} h(t) h(t - kT)\, dt = \delta_{k,0}
\]

Esto implica que no hay interferencia entre s\'imbolos (ISI) y la detecci\'on puede hacerse s\'imbolo por s\'imbolo.

\textbf{Deducci\'on de la normalizaci\'on:}

Para $k = 0$:
\[
\int_{\mathbb{R}} h(t) h(t)\, dt = \int_{\mathbb{R}} |h(t)|^2\, dt = 1
\]

Por lo tanto: $\boxed{\int_{\mathbb{R}} |h(t)|^2\, dt = 1}$

\textbf{Interpretaci\'on f\'isica:}
\begin{itemize}
    \item El filtro est\'a normalizado en energ\'ia unitaria
    \item Esta normalizaci\'on simplifica las expresiones del SNR
    \item Garantiza que la potencia del s\'imbolo dependa solo de la amplitud $A$ y los s\'imbolos $a_k, b_k$
\end{itemize}



\subsection*{Pregunta 7: Expresi\'on de $z_k$}

\textbf{Enunciado:} Calcular la expresi\'on de $z_k$ definida como la entrada del dispositivo de decisi\'on desde el cual se toma la decisi\'on sobre $d_k$.

\vspace{0.3cm}
\textbf{Soluci\'on:}

\textbf{An\'alisis de la rama I (In-phase):}

\begin{align*}
r_c(t) &= r(t) \cos(2\pi\nu_0 t) \\
&= \frac{A}{2}\sum_n a_n h(t-nT) + \frac{A}{2}\sum_n a_n h(t-nT)\cos(4\pi\nu_0 t) \\
&\quad + \text{t\'erminos en cuadratura} + n_c(t)
\end{align*}

Tras el filtro adaptado $g(t) = h(t_0 - t)$ y el muestreo en $t = t_0 + kT$:

\begin{align*}
a_k' &= (r_c * g)(t_0 + kT) \\
&= \frac{A}{2}\sum_n a_n \int_{\mathbb{R}} h(\tau) h(t_0 + kT - \tau - nT)\, d\tau + w_{c,k} \\
&= \frac{A}{2}\sum_n a_n \int_{\mathbb{R}} h(\tau) h(\tau + (k-n)T)\, d\tau + w_{c,k}
\end{align*}

Por TSWOC:
\[
a_k' = \frac{A}{2} a_k + w_{c,k}
\]

\textbf{An\'alisis de la rama Q (Quadrature):}

De forma similar:
\[
b_k' = \frac{A}{2} b_k + w_{s,k}
\]

\textbf{S\'imbolo complejo a la entrada del decisor:}

\[
\boxed{z_k = a_k' + ib_k' = \frac{A}{2}(a_k + ib_k) + (w_{c,k} + iw_{s,k}) = \frac{A}{2}d_k + w_k}
\]

Donde:
\begin{itemize}
    \item $d_k = a_k + ib_k$ es el s\'imbolo transmitido
    \item $w_k = w_{c,k} + iw_{s,k}$ es el ruido complejo gaussiano con varianza $\sigma^2 = \frac{N_0}{2}$ por componente
\end{itemize}

\textbf{Interpretaci\'on:}
\begin{itemize}
    \item La salida es una versi\'on escalada del s\'imbolo transmitido m\'as ruido aditivo
    \item El factor $\frac{A}{2}$ viene de la demodulaci\'on coherente
    \item No hay interferencia entre s\'imbolos gracias al TSWOC
\end{itemize}



\subsection*{Pregunta 8: Constelaci\'on recibida y regiones de decisi\'on}

\textbf{Enunciado:} A partir de la pregunta 7, trazar la constelaci\'on recibida asumiendo que el ruido es cero. Dibujar las regiones de decisi\'on binarias.

\vspace{0.3cm}
\textbf{Soluci\'on:}

\textbf{Constelaci\'on recibida (sin ruido):}

Con ruido cero, $z_k = \frac{A}{2}d_k$ donde $d_k \in \mathcal{D}$.

La constelaci\'on recibida es:
\[
\left\{\frac{A}{2}e^{i\frac{(2k+1)\pi}{8}} \mid k \in \{0,1,2,3,4,5,6,7\}\right\}
\]

\textbf{Representaci\'on gr\'afica con Gray mapping:}

\begin{center}
\textit{[Constelaci\'on 8-PSK con s\'imbolos en \'angulos m\'ultiplos de 45$^\circ$]}
\begin{itemize}
    \item $\frac{\pi}{8}$ (22.5$^\circ$): 000
    \item $\frac{3\pi}{8}$ (67.5$^\circ$): 001
    \item $\frac{5\pi}{8}$ (112.5$^\circ$): 011
    \item $\frac{7\pi}{8}$ (157.5$^\circ$): 010
    \item $\frac{9\pi}{8}$ (202.5$^\circ$): 110
    \item $\frac{11\pi}{8}$ (247.5$^\circ$): 111
    \item $\frac{13\pi}{8}$ (292.5$^\circ$): 101
    \item $\frac{15\pi}{8}$ (337.5$^\circ$): 100
\end{itemize}
\end{center}

\textbf{Regiones de decisi\'on binarias:}

Las fronteras de decisi\'on son las mediatrices perpendiculares entre s\'imbolos adyacentes, que pasan por el origen y est\'an en los \'angulos $0°, 45°, 90°, 135°, 180°, 225°, 270°, 315°$.

\textbf{Regla de decisi\'on ML:}

Decidir el s\'imbolo $\hat{d}_k$ m\'as cercano a $z_k$:
\[
\hat{d}_k = \arg\min_{d \in \mathcal{D}} |z_k - d|^2
\]

Equivalentemente, en t\'erminos del \'angulo:
\[
\text{Si } \arg(z_k) \in \left[\frac{(2k+1)\pi}{8} - \frac{\pi}{8}, \frac{(2k+1)\pi}{8} + \frac{\pi}{8}\right] \Rightarrow \hat{d}_k = e^{i\frac{(2k+1)\pi}{8}}
\]

\textbf{Decisi\'on binaria:}

Para cada bit del c\'odigo Gray, hay una frontera que separa los s\'imbolos con ese bit a 0 de los que lo tienen a 1.

\textbf{Ejemplo - Bit m\'as significativo (MSB):}
\begin{itemize}
    \item MSB = 0: s\'imbolos en el semiplano derecho (Re$(z_k) > 0$)
    \item MSB = 1: s\'imbolos en el semiplano izquierdo (Re$(z_k) < 0$)
\end{itemize}



\vspace{1cm}
\begin{center}
\rule{0.8\linewidth}{0.4pt}\\
\textit{Fin del examen}
\end{center}

\end{document}
